
\documentclass{article}
\usepackage{colortbl}
\usepackage{makecell}
\usepackage{multirow}
\usepackage{supertabular}

\begin{document}

\newcounter{utterance}

\centering \large Interaction Transcript for game `clean\_up', experiment `0\_easy\_3obj\_es', episode 4 with gemma{-}4{-}26b{-}a4b{-}it.
\vspace{24pt}

{ \footnotesize  \setcounter{utterance}{1}
\setlength{\tabcolsep}{0pt}
\begin{supertabular}{c@{$\;$}|p{.15\linewidth}@{}p{.15\linewidth}p{.15\linewidth}p{.15\linewidth}p{.15\linewidth}p{.15\linewidth}}
    \# & $\;$A & \multicolumn{4}{c}{Game Master} & $\;\:$B\\
    \hline

    \theutterance \stepcounter{utterance}  
    & & \multicolumn{4}{p{0.6\linewidth}}{
        \cellcolor[rgb]{0.9,0.9,0.9}{
            \makecell[{{p{\linewidth}}}]{
                \texttt{\tiny{[P1$\langle$GM]}}
                \texttt{Soy tu maestro de juego, y estás jugando un juego colaborativo con la siguiente cuadrícula como tablero de juego:} \\
\\ 
\texttt{```} \\
\texttt{ 1234567} \\
\texttt{╔═╤═════╗} \\
\texttt{║◌│◌◌◌◌L║ 1} \\
\texttt{║◌│◌◌◌◌◌║ 2} \\
\texttt{║◌│◌◌◌◌◌║ 3} \\
\texttt{║◌│◌◌◌◌◌║ 4} \\
\texttt{╟─┼───┐C║ 5} \\
\texttt{║◌│◌◌◌├─╢ 6} \\
\texttt{║P│◌◌◌│◌║ 7} \\
\texttt{╚═╧═══╧═╝} \\
\\ 
\texttt{```} \\
\\ 
\texttt{* El borde superior muestra coordenadas x que aumentan hacia la derecha, y el borde derecho coordenadas y que aumentan hacia abajo.} \\
\texttt{* Los siguientes objetos están colocados aleatoriamente en tu cuadrícula: 'C', 'L', 'P'.} \\
\\ 
\\ 
\texttt{El otro jugador ve una variación del tablero de juego, donde los objetos están colocados en diferentes ubicaciones aleatorias. No puedes ver el tablero del otro jugador, y él no puede ver el tuyo.} \\
\\ 
\\ 
\texttt{**Objetivo:**} \\
\\ 
\texttt{Ambos jugadores necesitan mover los objetos en sus respectivos fondos de manera que los objetos idénticos terminen en las mismas coordenadas. Tienes que comunicarte con el otro jugador para acordar una configuración objetivo común.} \\
\\ 
\\ 
\texttt{**Reglas:**} \\
\\ 
\texttt{* En cada turno, puedes enviar exactamente uno de los dos comandos siguientes:} \\
\texttt{ 1. `DECIR: <MENSAJE>`: para enviar un mensaje (todo hasta el siguiente salto de línea) al otro jugador. Se lo reenviaré a tu compañero.} \\
\texttt{2. `MOVERR: <OBJETO>, (<X>, <Y>)`: para mover un objeto a una nueva posición, donde `<X>` es la columna y `<Y>` es la fila. Te informaré si tu movimiento fue exitoso o no.} \\
\texttt{* Si no sigues el formato, o envías varios comandos a la vez, tendré que penalizarte.} \\
\texttt{* Si ambos jugadores acumulan más de 8 penalizaciones, ambos pierden el juego.} \\
\texttt{* ¡Es vital que te comuniques con el otro jugador respecto a tu estado objetivo! ¡La *única* forma de transmitir tu estrategia al otro jugador es usando el comando `DECIR: <MENSAJE>`!} \\
\\ 
\\ 
\texttt{**Mover objetos**} \\
\\ 
\texttt{* Solo puedes mover objetos a celdas dentro de los límites de la cuadrícula. La celda de destino debe estar vacía, es decir, solo debe contener el símbolo '◌'.} \\
\texttt{* Si intentas mover un objeto a un lugar que no está vacío, o intentas moverlo fuera de la cuadrícula, tendré que penalizarte. Tendrás otro intento.} \\
\texttt{* ¡Antes de hacer un movimiento, verifica dos veces que el lugar de destino esté vacío, y no contenga ninguna letra, marco o línea!} \\
\\ 
\\ 
\texttt{**Fin del juego**} \\
\\ 
\texttt{Si crees que alcanzaste el objetivo de alinear todos los objetos, puedes pedirle al otro jugador que termine el juego enviando `DECIR: ¿terminado?`. Si el otro jugador te pide que termines el juego, y respondes con `DECIR: ¡terminado!`, el juego terminará.} \\
\\ 
\texttt{Ambos jugadores ganan si el juego termina dentro de 12 rondas, donde una ronda se define como dos jugadores enviando cada uno un comando válido.} \\
\\ 
\texttt{**Puntuación:**} \\
\\ 
\texttt{Cuanto más cerca estén los objetos idénticos en ambos tableros de juego, más puntos obtendrás. Las penalizaciones reducen tus puntos. ¿Puedes batir el récord?} \\
\\ 
\\ 
\texttt{¡Por favor, envía un mensaje al otro jugador para iniciar el juego!} \\
            }
        }
    }
    & \\ \\

    \theutterance \stepcounter{utterance}  
    & \multicolumn{4}{p{0.6\linewidth}}{
        \cellcolor[rgb]{0.9,0.9,0.9}{
            \makecell[{{p{\linewidth}}}]{
                \texttt{\tiny{[P1$\rangle$GM]}}
                \texttt{DECIR: Hola compañero. Para empezar, voy a decirte mis posiciones actuales para que podamos coordinarnos. En mi tablero, los objetos están en: L en (6,1), C en (6,5) y P en (1,7). ¿Podrías decirme dónde están los tuyos?} \\
            }
        }
    }
    & & \\ \\

    \theutterance \stepcounter{utterance}  
    & & & \multicolumn{2}{p{0.3\linewidth}}{
        \cellcolor[rgb]{0.9,0.9,0.9}{
            \makecell[{{p{\linewidth}}}]{
                \texttt{\tiny{[GM$|$GM]}}
                \texttt{DECIR: Hola compañero. Para empezar, voy a decirte mis posiciones actuales para que podamos coordinarnos. En mi tablero, los objetos están en: L en (6,1), C en (6,5) y P en (1,7). ¿Podrías decirme dónde están los tuyos?} \\
            }
        }
    }
    & & \\ \\

    \theutterance \stepcounter{utterance}  
    & & \multicolumn{4}{p{0.6\linewidth}}{
        \cellcolor[rgb]{0.9,0.9,0.9}{
            \makecell[{{p{\linewidth}}}]{
                \texttt{\tiny{[GM$\rangle$P2]}}
                \texttt{Soy tu maestro de juego, y estás jugando un juego colaborativo con la siguiente cuadrícula como tablero de juego:} \\
\\ 
\texttt{```} \\
\texttt{ 1234567} \\
\texttt{╔═╤═════╗} \\
\texttt{║◌│◌◌◌◌◌║ 1} \\
\texttt{║◌│◌◌PL◌║ 2} \\
\texttt{║◌│◌◌◌◌◌║ 3} \\
\texttt{║◌│◌◌◌◌◌║ 4} \\
\texttt{╟─┼───┐◌║ 5} \\
\texttt{║◌│◌◌◌├─╢ 6} \\
\texttt{║◌│◌C◌│◌║ 7} \\
\texttt{╚═╧═══╧═╝} \\
\\ 
\texttt{```} \\
\\ 
\texttt{* El borde superior muestra coordenadas x que aumentan hacia la derecha, y el borde derecho coordenadas y que aumentan hacia abajo.} \\
\texttt{* Los siguientes objetos están colocados aleatoriamente en tu cuadrícula: 'C', 'L', 'P'.} \\
\\ 
\\ 
\texttt{El otro jugador ve una variación del tablero de juego, donde los objetos están colocados en diferentes ubicaciones aleatorias. No puedes ver el tablero del otro jugador, y él no puede ver el tuyo.} \\
\\ 
\\ 
\texttt{**Objetivo:**} \\
\\ 
\texttt{Ambos jugadores necesitan mover los objetos en sus respectivos fondos de manera que los objetos idénticos terminen en las mismas coordenadas. Tienes que comunicarte con el otro jugador para acordar una configuración objetivo común.} \\
\\ 
\\ 
\texttt{**Reglas:**} \\
\\ 
\texttt{* En cada turno, puedes enviar exactamente uno de los dos comandos siguientes:} \\
\texttt{ 1. `DECIR: <MENSAJE>`: para enviar un mensaje (todo hasta el siguiente salto de línea) al otro jugador. Se lo reenviaré a tu compañero.} \\
\texttt{2. `MOVERR: <OBJETO>, (<X>, <Y>)`: para mover un objeto a una nueva posición, donde `<X>` es la columna y `<Y>` es la fila. Te informaré si tu movimiento fue exitoso o no.} \\
\texttt{* Si no sigues el formato, o envías varios comandos a la vez, tendré que penalizarte.} \\
\texttt{* Si ambos jugadores acumulan más de 8 penalizaciones, ambos pierden el juego.} \\
\texttt{* ¡Es vital que te comuniques con el otro jugador respecto a tu estado objetivo! ¡La *única* forma de transmitir tu estrategia al otro jugador es usando el comando `DECIR: <MENSAJE>`!} \\
\\ 
\\ 
\texttt{**Mover objetos**} \\
\\ 
\texttt{* Solo puedes mover objetos a celdas dentro de los límites de la cuadrícula. La celda de destino debe estar vacía, es decir, solo debe contener el símbolo '◌'.} \\
\texttt{* Si intentas mover un objeto a un lugar que no está vacío, o intentas moverlo fuera de la cuadrícula, tendré que penalizarte. Tendrás otro intento.} \\
\texttt{* ¡Antes de hacer un movimiento, verifica dos veces que el lugar de destino esté vacío, y no contenga ninguna letra, marco o línea!} \\
\\ 
\\ 
\texttt{**Fin del juego**} \\
\\ 
\texttt{Si crees que alcanzaste el objetivo de alinear todos los objetos, puedes pedirle al otro jugador que termine el juego enviando `DECIR: ¿terminado?`. Si el otro jugador te pide que termines el juego, y respondes con `DECIR: ¡terminado!`, el juego terminará.} \\
\\ 
\texttt{Ambos jugadores ganan si el juego termina dentro de 12 rondas, donde una ronda se define como dos jugadores enviando cada uno un comando válido.} \\
\\ 
\texttt{**Puntuación:**} \\
\\ 
\texttt{Cuanto más cerca estén los objetos idénticos en ambos tableros de juego, más puntos obtendrás. Las penalizaciones reducen tus puntos. ¿Puedes batir el récord?} \\
\\ 
\texttt{El otro jugador inició el juego enviando este mensaje:} \\
\\ 
\texttt{"Hola compañero. Para empezar, voy a decirte mis posiciones actuales para que podamos coordinarnos. En mi tablero, los objetos están en: L en (6,1), C en (6,5) y P en (1,7). ¿Podrías decirme dónde están los tuyos?"} \\
\\ 
\texttt{¿Cuál es tu primer comando?} \\
            }
        }
    }
    & \\ \\

    \theutterance \stepcounter{utterance}  
    & & & \multicolumn{4}{p{0.6\linewidth}}{
        \cellcolor[rgb]{0.9,0.9,0.9}{
            \makecell[{{p{\linewidth}}}]{
                \texttt{\tiny{[GM$\langle$P2]}}
                \texttt{DECIR: Hola. Mis posiciones actuales son: P en (5,2), L en (6,2) y C en (4,7). ¿Qué configuración objetivo te parece mejor para que todos coincidan?} \\
            }
        }
    }
     \\ \\

    \theutterance \stepcounter{utterance}  
    & & & \multicolumn{2}{p{0.3\linewidth}}{
        \cellcolor[rgb]{0.9,0.9,0.9}{
            \makecell[{{p{\linewidth}}}]{
                \texttt{\tiny{[GM$|$GM]}}
                \texttt{DECIR: Hola. Mis posiciones actuales son: P en (5,2), L en (6,2) y C en (4,7). ¿Qué configuración objetivo te parece mejor para que todos coincidan?} \\
            }
        }
    }
    & & \\ \\

    \theutterance \stepcounter{utterance}  
    & & \multicolumn{4}{p{0.6\linewidth}}{
        \cellcolor[rgb]{0.9,0.9,0.9}{
            \makecell[{{p{\linewidth}}}]{
                \texttt{\tiny{[P1$\langle$GM]}}
                \texttt{Tu mensaje ha sido transmitido al otro jugador.} \\
\\ 
\\ 
\texttt{Actualmente estás jugando la ronda 1 de un máximo de 12.} \\
\texttt{Han acumulado colectivamente 0 de 8 penalizaciones.} \\
\texttt{El otro jugador envió este mensaje:} \\
\\ 
\\ 
\texttt{"Hola. Mis posiciones actuales son: P en (5,2), L en (6,2) y C en (4,7). ¿Qué configuración objetivo te parece mejor para que todos coincidan?"} \\
\\ 
\\ 
\texttt{¿Cuál es tu próximo comando?} \\
            }
        }
    }
    & \\ \\

    \theutterance \stepcounter{utterance}  
    & \multicolumn{4}{p{0.6\linewidth}}{
        \cellcolor[rgb]{0.9,0.9,0.9}{
            \makecell[{{p{\linewidth}}}]{
                \texttt{\tiny{[P1$\rangle$GM]}}
                \texttt{DECIR: Gracias por la información. Mis posiciones son L(6,1), C(6,5) y P(1,7). Tus posiciones son P(5,2), L(6,2) y C(4,7). Propongo que movamos todos los objetos a una configuración central para minimizar los movimientos. ¿Qué te parece si fijamos el objetivo en: L en (6,2), C en (5,5) y P en (5,2)?} \\
            }
        }
    }
    & & \\ \\

    \theutterance \stepcounter{utterance}  
    & & & \multicolumn{2}{p{0.3\linewidth}}{
        \cellcolor[rgb]{0.9,0.9,0.9}{
            \makecell[{{p{\linewidth}}}]{
                \texttt{\tiny{[GM$|$GM]}}
                \texttt{DECIR: Gracias por la información. Mis posiciones son L(6,1), C(6,5) y P(1,7). Tus posiciones son P(5,2), L(6,2) y C(4,7). Propongo que movamos todos los objetos a una configuración central para minimizar los movimientos. ¿Qué te parece si fijamos el objetivo en: L en (6,2), C en (5,5) y P en (5,2)?} \\
            }
        }
    }
    & & \\ \\

    \theutterance \stepcounter{utterance}  
    & & \multicolumn{4}{p{0.6\linewidth}}{
        \cellcolor[rgb]{0.9,0.9,0.9}{
            \makecell[{{p{\linewidth}}}]{
                \texttt{\tiny{[GM$\rangle$P2]}}
                \texttt{Tu mensaje ha sido transmitido al otro jugador.} \\
\\ 
\\ 
\texttt{Actualmente estás jugando la ronda 2 de un máximo de 12.} \\
\texttt{Han acumulado colectivamente 0 de 8 penalizaciones.} \\
\texttt{El otro jugador envió este mensaje:} \\
\\ 
\\ 
\texttt{"Gracias por la información. Mis posiciones son L(6,1), C(6,5) y P(1,7). Tus posiciones son P(5,2), L(6,2) y C(4,7). Propongo que movamos todos los objetos a una configuración central para minimizar los movimientos. ¿Qué te parece si fijamos el objetivo en: L en (6,2), C en (5,5) y P en (5,2)?"} \\
\\ 
\\ 
\texttt{¿Cuál es tu próximo comando?} \\
            }
        }
    }
    & \\ \\

    \theutterance \stepcounter{utterance}  
    & & & \multicolumn{4}{p{0.6\linewidth}}{
        \cellcolor[rgb]{0.9,0.9,0.9}{
            \makecell[{{p{\linewidth}}}]{
                \texttt{\tiny{[GM$\langle$P2]}}
                \texttt{DECIR: Me parece una excelente propuesta. Es una configuración muy equilibrada. Empezaré moviendo mi objeto C hacia el objetivo. MOVERR: C, (5,5)} \\
            }
        }
    }
     \\ \\

    \theutterance \stepcounter{utterance}  
    & & & \multicolumn{2}{p{0.3\linewidth}}{
        \cellcolor[rgb]{0.9,0.9,0.9}{
            \makecell[{{p{\linewidth}}}]{
                \texttt{\tiny{[GM$|$GM]}}
                \texttt{DECIR: Me parece una excelente propuesta. Es una configuración muy equilibrada. Empezaré moviendo mi objeto C hacia el objetivo. MOVERR: C, (5,5)} \\
            }
        }
    }
    & & \\ \\

    \theutterance \stepcounter{utterance}  
    & & \multicolumn{4}{p{0.6\linewidth}}{
        \cellcolor[rgb]{0.9,0.9,0.9}{
            \makecell[{{p{\linewidth}}}]{
                \texttt{\tiny{[P1$\langle$GM]}}
                \texttt{Tu mensaje ha sido transmitido al otro jugador.} \\
\\ 
\\ 
\texttt{Actualmente estás jugando la ronda 2 de un máximo de 12.} \\
\texttt{Han acumulado colectivamente 0 de 8 penalizaciones.} \\
\texttt{El otro jugador envió este mensaje:} \\
\\ 
\\ 
\texttt{"Me parece una excelente propuesta. Es una configuración muy equilibrada. Empezaré moviendo mi objeto C hacia el objetivo. MOVERR: C, (5,5)"} \\
\\ 
\\ 
\texttt{¿Cuál es tu próximo comando?} \\
            }
        }
    }
    & \\ \\

    \theutterance \stepcounter{utterance}  
    & \multicolumn{4}{p{0.6\linewidth}}{
        \cellcolor[rgb]{0.9,0.9,0.9}{
            \makecell[{{p{\linewidth}}}]{
                \texttt{\tiny{[P1$\rangle$GM]}}
                \texttt{MOVERR: L, (6,2)} \\
            }
        }
    }
    & & \\ \\

    \theutterance \stepcounter{utterance}  
    & & & \multicolumn{2}{p{0.3\linewidth}}{
        \cellcolor[rgb]{0.9,0.9,0.9}{
            \makecell[{{p{\linewidth}}}]{
                \texttt{\tiny{[GM$|$GM]}}
                \texttt{MOVERR: L, (6,2)} \\
            }
        }
    }
    & & \\ \\

    \theutterance \stepcounter{utterance}  
    & & & \multicolumn{2}{p{0.3\linewidth}}{
        \cellcolor[rgb]{0.9,0.9,0.9}{
            \makecell[{{p{\linewidth}}}]{
                \texttt{\tiny{[GM$|$GM]}}
                \texttt{Invalid response format} \\
            }
        }
    }
    & & \\ \\

    \theutterance \stepcounter{utterance}  
    & & \multicolumn{4}{p{0.6\linewidth}}{
        \cellcolor[rgb]{0.9,0.9,0.9}{
            \makecell[{{p{\linewidth}}}]{
                \texttt{\tiny{[P1$\langle$GM]}}
                \texttt{Penalización: ¡Tu mensaje no tiene el formato esperado!} \\
\\ 
\texttt{Asegúrate de que tu respuesta solo contenga `DECIR: <MENSAJE>` o `MOVERR: <OBJETO>, (<X>, <Y>)`, ¡y nada más!} \\
\\ 
\texttt{Han acumulado colectivamente 1 de 8 penalizaciones.} \\
\texttt{¡Por favor, inténtalo de nuevo!} \\
            }
        }
    }
    & \\ \\

    \theutterance \stepcounter{utterance}  
    & \multicolumn{4}{p{0.6\linewidth}}{
        \cellcolor[rgb]{0.9,0.9,0.9}{
            \makecell[{{p{\linewidth}}}]{
                \texttt{\tiny{[P1$\rangle$GM]}}
                \texttt{MOVERR: L, (6,2)} \\
            }
        }
    }
    & & \\ \\

    \theutterance \stepcounter{utterance}  
    & & & \multicolumn{2}{p{0.3\linewidth}}{
        \cellcolor[rgb]{0.9,0.9,0.9}{
            \makecell[{{p{\linewidth}}}]{
                \texttt{\tiny{[GM$|$GM]}}
                \texttt{MOVERR: L, (6,2)} \\
            }
        }
    }
    & & \\ \\

    \theutterance \stepcounter{utterance}  
    & & & \multicolumn{2}{p{0.3\linewidth}}{
        \cellcolor[rgb]{0.9,0.9,0.9}{
            \makecell[{{p{\linewidth}}}]{
                \texttt{\tiny{[GM$|$GM]}}
                \texttt{Invalid response format} \\
            }
        }
    }
    & & \\ \\

    \theutterance \stepcounter{utterance}  
    & & \multicolumn{4}{p{0.6\linewidth}}{
        \cellcolor[rgb]{0.9,0.9,0.9}{
            \makecell[{{p{\linewidth}}}]{
                \texttt{\tiny{[P1$\langle$GM]}}
                \texttt{Penalización: ¡Tu mensaje no tiene el formato esperado!} \\
\\ 
\texttt{Asegúrate de que tu respuesta solo contenga `DECIR: <MENSAJE>` o `MOVERR: <OBJETO>, (<X>, <Y>)`, ¡y nada más!} \\
\\ 
\texttt{Han acumulado colectivamente 2 de 8 penalizaciones.} \\
\texttt{¡Por favor, inténtalo de nuevo!} \\
            }
        }
    }
    & \\ \\

    \theutterance \stepcounter{utterance}  
    & \multicolumn{4}{p{0.6\linewidth}}{
        \cellcolor[rgb]{0.9,0.9,0.9}{
            \makecell[{{p{\linewidth}}}]{
                \texttt{\tiny{[P1$\rangle$GM]}}
                \texttt{MOVERR: L, (6,2)} \\
            }
        }
    }
    & & \\ \\

    \theutterance \stepcounter{utterance}  
    & & & \multicolumn{2}{p{0.3\linewidth}}{
        \cellcolor[rgb]{0.9,0.9,0.9}{
            \makecell[{{p{\linewidth}}}]{
                \texttt{\tiny{[GM$|$GM]}}
                \texttt{MOVERR: L, (6,2)} \\
            }
        }
    }
    & & \\ \\

    \theutterance \stepcounter{utterance}  
    & & & \multicolumn{2}{p{0.3\linewidth}}{
        \cellcolor[rgb]{0.9,0.9,0.9}{
            \makecell[{{p{\linewidth}}}]{
                \texttt{\tiny{[GM$|$GM]}}
                \texttt{Invalid response format} \\
            }
        }
    }
    & & \\ \\

    \theutterance \stepcounter{utterance}  
    & & \multicolumn{4}{p{0.6\linewidth}}{
        \cellcolor[rgb]{0.9,0.9,0.9}{
            \makecell[{{p{\linewidth}}}]{
                \texttt{\tiny{[P1$\langle$GM]}}
                \texttt{Penalización: ¡Tu mensaje no tiene el formato esperado!} \\
\\ 
\texttt{Asegúrate de que tu respuesta solo contenga `DECIR: <MENSAJE>` o `MOVERR: <OBJETO>, (<X>, <Y>)`, ¡y nada más!} \\
\\ 
\texttt{Han acumulado colectivamente 3 de 8 penalizaciones.} \\
\texttt{¡Por favor, inténtalo de nuevo!} \\
            }
        }
    }
    & \\ \\

    \theutterance \stepcounter{utterance}  
    & \multicolumn{4}{p{0.6\linewidth}}{
        \cellcolor[rgb]{0.9,0.9,0.9}{
            \makecell[{{p{\linewidth}}}]{
                \texttt{\tiny{[P1$\rangle$GM]}}
                \texttt{MOVERR: L, (6,2)} \\
            }
        }
    }
    & & \\ \\

    \theutterance \stepcounter{utterance}  
    & & & \multicolumn{2}{p{0.3\linewidth}}{
        \cellcolor[rgb]{0.9,0.9,0.9}{
            \makecell[{{p{\linewidth}}}]{
                \texttt{\tiny{[GM$|$GM]}}
                \texttt{MOVERR: L, (6,2)} \\
            }
        }
    }
    & & \\ \\

    \theutterance \stepcounter{utterance}  
    & & & \multicolumn{2}{p{0.3\linewidth}}{
        \cellcolor[rgb]{0.9,0.9,0.9}{
            \makecell[{{p{\linewidth}}}]{
                \texttt{\tiny{[GM$|$GM]}}
                \texttt{Invalid response format} \\
            }
        }
    }
    & & \\ \\

    \theutterance \stepcounter{utterance}  
    & & \multicolumn{4}{p{0.6\linewidth}}{
        \cellcolor[rgb]{0.9,0.9,0.9}{
            \makecell[{{p{\linewidth}}}]{
                \texttt{\tiny{[P1$\langle$GM]}}
                \texttt{Penalización: ¡Tu mensaje no tiene el formato esperado!} \\
\\ 
\texttt{Asegúrate de que tu respuesta solo contenga `DECIR: <MENSAJE>` o `MOVERR: <OBJETO>, (<X>, <Y>)`, ¡y nada más!} \\
\\ 
\texttt{Han acumulado colectivamente 4 de 8 penalizaciones.} \\
\texttt{¡Por favor, inténtalo de nuevo!} \\
            }
        }
    }
    & \\ \\

    \theutterance \stepcounter{utterance}  
    & \multicolumn{4}{p{0.6\linewidth}}{
        \cellcolor[rgb]{0.9,0.9,0.9}{
            \makecell[{{p{\linewidth}}}]{
                \texttt{\tiny{[P1$\rangle$GM]}}
                \texttt{MOVERR: L, (6,2)} \\
            }
        }
    }
    & & \\ \\

    \theutterance \stepcounter{utterance}  
    & & & \multicolumn{2}{p{0.3\linewidth}}{
        \cellcolor[rgb]{0.9,0.9,0.9}{
            \makecell[{{p{\linewidth}}}]{
                \texttt{\tiny{[GM$|$GM]}}
                \texttt{MOVERR: L, (6,2)} \\
            }
        }
    }
    & & \\ \\

    \theutterance \stepcounter{utterance}  
    & & & \multicolumn{2}{p{0.3\linewidth}}{
        \cellcolor[rgb]{0.9,0.9,0.9}{
            \makecell[{{p{\linewidth}}}]{
                \texttt{\tiny{[GM$|$GM]}}
                \texttt{Invalid response format} \\
            }
        }
    }
    & & \\ \\

    \theutterance \stepcounter{utterance}  
    & & \multicolumn{4}{p{0.6\linewidth}}{
        \cellcolor[rgb]{0.9,0.9,0.9}{
            \makecell[{{p{\linewidth}}}]{
                \texttt{\tiny{[P1$\langle$GM]}}
                \texttt{Penalización: ¡Tu mensaje no tiene el formato esperado!} \\
\\ 
\texttt{Asegúrate de que tu respuesta solo contenga `DECIR: <MENSAJE>` o `MOVERR: <OBJETO>, (<X>, <Y>)`, ¡y nada más!} \\
\\ 
\texttt{Han acumulado colectivamente 5 de 8 penalizaciones.} \\
\texttt{¡Por favor, inténtalo de nuevo!} \\
            }
        }
    }
    & \\ \\

    \theutterance \stepcounter{utterance}  
    & \multicolumn{4}{p{0.6\linewidth}}{
        \cellcolor[rgb]{0.9,0.9,0.9}{
            \makecell[{{p{\linewidth}}}]{
                \texttt{\tiny{[P1$\rangle$GM]}}
                \texttt{MOVERR: L, (6,2)} \\
            }
        }
    }
    & & \\ \\

    \theutterance \stepcounter{utterance}  
    & & & \multicolumn{2}{p{0.3\linewidth}}{
        \cellcolor[rgb]{0.9,0.9,0.9}{
            \makecell[{{p{\linewidth}}}]{
                \texttt{\tiny{[GM$|$GM]}}
                \texttt{MOVERR: L, (6,2)} \\
            }
        }
    }
    & & \\ \\

    \theutterance \stepcounter{utterance}  
    & & & \multicolumn{2}{p{0.3\linewidth}}{
        \cellcolor[rgb]{0.9,0.9,0.9}{
            \makecell[{{p{\linewidth}}}]{
                \texttt{\tiny{[GM$|$GM]}}
                \texttt{Invalid response format} \\
            }
        }
    }
    & & \\ \\

    \theutterance \stepcounter{utterance}  
    & & \multicolumn{4}{p{0.6\linewidth}}{
        \cellcolor[rgb]{0.9,0.9,0.9}{
            \makecell[{{p{\linewidth}}}]{
                \texttt{\tiny{[P1$\langle$GM]}}
                \texttt{Penalización: ¡Tu mensaje no tiene el formato esperado!} \\
\\ 
\texttt{Asegúrate de que tu respuesta solo contenga `DECIR: <MENSAJE>` o `MOVERR: <OBJETO>, (<X>, <Y>)`, ¡y nada más!} \\
\\ 
\texttt{Han acumulado colectivamente 6 de 8 penalizaciones.} \\
\texttt{¡Por favor, inténtalo de nuevo!} \\
            }
        }
    }
    & \\ \\

    \theutterance \stepcounter{utterance}  
    & \multicolumn{4}{p{0.6\linewidth}}{
        \cellcolor[rgb]{0.9,0.9,0.9}{
            \makecell[{{p{\linewidth}}}]{
                \texttt{\tiny{[P1$\rangle$GM]}}
                \texttt{MOVERR: L, (6,2)} \\
            }
        }
    }
    & & \\ \\

    \theutterance \stepcounter{utterance}  
    & & & \multicolumn{2}{p{0.3\linewidth}}{
        \cellcolor[rgb]{0.9,0.9,0.9}{
            \makecell[{{p{\linewidth}}}]{
                \texttt{\tiny{[GM$|$GM]}}
                \texttt{MOVERR: L, (6,2)} \\
            }
        }
    }
    & & \\ \\

    \theutterance \stepcounter{utterance}  
    & & & \multicolumn{2}{p{0.3\linewidth}}{
        \cellcolor[rgb]{0.9,0.9,0.9}{
            \makecell[{{p{\linewidth}}}]{
                \texttt{\tiny{[GM$|$GM]}}
                \texttt{Invalid response format} \\
            }
        }
    }
    & & \\ \\

    \theutterance \stepcounter{utterance}  
    & & \multicolumn{4}{p{0.6\linewidth}}{
        \cellcolor[rgb]{0.9,0.9,0.9}{
            \makecell[{{p{\linewidth}}}]{
                \texttt{\tiny{[P1$\langle$GM]}}
                \texttt{Penalización: ¡Tu mensaje no tiene el formato esperado!} \\
\\ 
\texttt{Asegúrate de que tu respuesta solo contenga `DECIR: <MENSAJE>` o `MOVERR: <OBJETO>, (<X>, <Y>)`, ¡y nada más!} \\
\\ 
\texttt{Han acumulado colectivamente 7 de 8 penalizaciones.} \\
\texttt{¡Por favor, inténtalo de nuevo!} \\
            }
        }
    }
    & \\ \\

    \theutterance \stepcounter{utterance}  
    & \multicolumn{4}{p{0.6\linewidth}}{
        \cellcolor[rgb]{0.9,0.9,0.9}{
            \makecell[{{p{\linewidth}}}]{
                \texttt{\tiny{[P1$\rangle$GM]}}
                \texttt{MOVERR: L, (6,2)} \\
            }
        }
    }
    & & \\ \\

    \theutterance \stepcounter{utterance}  
    & & & \multicolumn{2}{p{0.3\linewidth}}{
        \cellcolor[rgb]{0.9,0.9,0.9}{
            \makecell[{{p{\linewidth}}}]{
                \texttt{\tiny{[GM$|$GM]}}
                \texttt{MOVERR: L, (6,2)} \\
            }
        }
    }
    & & \\ \\

    \theutterance \stepcounter{utterance}  
    & & & \multicolumn{2}{p{0.3\linewidth}}{
        \cellcolor[rgb]{0.9,0.9,0.9}{
            \makecell[{{p{\linewidth}}}]{
                \texttt{\tiny{[GM$|$GM]}}
                \texttt{Invalid response format} \\
            }
        }
    }
    & & \\ \\

    \theutterance \stepcounter{utterance}  
    & & \multicolumn{4}{p{0.6\linewidth}}{
        \cellcolor[rgb]{0.9,0.9,0.9}{
            \makecell[{{p{\linewidth}}}]{
                \texttt{\tiny{[P1$\langle$GM]}}
                \texttt{Penalización: ¡Tu mensaje no tiene el formato esperado!} \\
\\ 
\texttt{Asegúrate de que tu respuesta solo contenga `DECIR: <MENSAJE>` o `MOVERR: <OBJETO>, (<X>, <Y>)`, ¡y nada más!} \\
\\ 
\texttt{Han acumulado colectivamente 8 de 8 penalizaciones.} \\
\texttt{¡Por favor, inténtalo de nuevo!} \\
            }
        }
    }
    & \\ \\

    \theutterance \stepcounter{utterance}  
    & \multicolumn{4}{p{0.6\linewidth}}{
        \cellcolor[rgb]{0.9,0.9,0.9}{
            \makecell[{{p{\linewidth}}}]{
                \texttt{\tiny{[P1$\rangle$GM]}}
                \texttt{MOVERR: L, (6,2)} \\
            }
        }
    }
    & & \\ \\

    \theutterance \stepcounter{utterance}  
    & & & \multicolumn{2}{p{0.3\linewidth}}{
        \cellcolor[rgb]{0.9,0.9,0.9}{
            \makecell[{{p{\linewidth}}}]{
                \texttt{\tiny{[GM$|$GM]}}
                \texttt{MOVERR: L, (6,2)} \\
            }
        }
    }
    & & \\ \\

    \theutterance \stepcounter{utterance}  
    & & & \multicolumn{2}{p{0.3\linewidth}}{
        \cellcolor[rgb]{0.9,0.9,0.9}{
            \makecell[{{p{\linewidth}}}]{
                \texttt{\tiny{[GM$|$GM]}}
                \texttt{Invalid response format} \\
            }
        }
    }
    & & \\ \\

    \theutterance \stepcounter{utterance}  
    & & & \multicolumn{2}{p{0.3\linewidth}}{
        \cellcolor[rgb]{0.9,0.9,0.9}{
            \makecell[{{p{\linewidth}}}]{
                \texttt{\tiny{[GM$|$GM]}}
                \texttt{Maximum number of penalties exceeded} \\
            }
        }
    }
    & & \\ \\

    \theutterance \stepcounter{utterance}  
    & & & \multicolumn{2}{p{0.3\linewidth}}{
        \cellcolor[rgb]{0.9,0.9,0.9}{
            \makecell[{{p{\linewidth}}}]{
                \texttt{\tiny{[GM$|$GM]}}
                \texttt{Initial states:} \\
\texttt{```} \\
\texttt{Player 1:} \\
\texttt{ 1234567} \\
\texttt{╔═╤═════╗} \\
\texttt{║◌│◌◌◌◌L║ 1} \\
\texttt{║◌│◌◌◌◌◌║ 2} \\
\texttt{║◌│◌◌◌◌◌║ 3} \\
\texttt{║◌│◌◌◌◌◌║ 4} \\
\texttt{╟─┼───┐C║ 5} \\
\texttt{║◌│◌◌◌├─╢ 6} \\
\texttt{║P│◌◌◌│◌║ 7} \\
\texttt{╚═╧═══╧═╝} \\
\\ 
\\ 
\texttt{Player 2:} \\
\texttt{ 1234567} \\
\texttt{╔═╤═════╗} \\
\texttt{║◌│◌◌◌◌◌║ 1} \\
\texttt{║◌│◌◌PL◌║ 2} \\
\texttt{║◌│◌◌◌◌◌║ 3} \\
\texttt{║◌│◌◌◌◌◌║ 4} \\
\texttt{╟─┼───┐◌║ 5} \\
\texttt{║◌│◌◌◌├─╢ 6} \\
\texttt{║◌│◌C◌│◌║ 7} \\
\texttt{╚═╧═══╧═╝} \\
\\ 
\texttt{```} \\
            }
        }
    }
    & & \\ \\

    \theutterance \stepcounter{utterance}  
    & & & \multicolumn{2}{p{0.3\linewidth}}{
        \cellcolor[rgb]{0.9,0.9,0.9}{
            \makecell[{{p{\linewidth}}}]{
                \texttt{\tiny{[GM$|$GM]}}
                \texttt{End states:} \\
\texttt{```} \\
\texttt{Player 1:} \\
\texttt{ 1234567} \\
\texttt{╔═╤═════╗} \\
\texttt{║◌│◌◌◌◌L║ 1} \\
\texttt{║◌│◌◌◌◌◌║ 2} \\
\texttt{║◌│◌◌◌◌◌║ 3} \\
\texttt{║◌│◌◌◌◌◌║ 4} \\
\texttt{╟─┼───┐C║ 5} \\
\texttt{║◌│◌◌◌├─╢ 6} \\
\texttt{║P│◌◌◌│◌║ 7} \\
\texttt{╚═╧═══╧═╝} \\
\\ 
\\ 
\texttt{Player 2:} \\
\texttt{ 1234567} \\
\texttt{╔═╤═════╗} \\
\texttt{║◌│◌◌◌◌◌║ 1} \\
\texttt{║◌│◌◌PL◌║ 2} \\
\texttt{║◌│◌◌◌◌◌║ 3} \\
\texttt{║◌│◌◌◌◌◌║ 4} \\
\texttt{╟─┼───┐◌║ 5} \\
\texttt{║◌│◌◌◌├─╢ 6} \\
\texttt{║◌│◌C◌│◌║ 7} \\
\texttt{╚═╧═══╧═╝} \\
\\ 
\texttt{```} \\
            }
        }
    }
    & & \\ \\

    \theutterance \stepcounter{utterance}  
    & & & \multicolumn{2}{p{0.3\linewidth}}{
        \cellcolor[rgb]{0.9,0.9,0.9}{
            \makecell[{{p{\linewidth}}}]{
                \texttt{\tiny{[GM$|$GM]}}
                \texttt{* success: False} \\
\texttt{* lose: False} \\
\texttt{* aborted: True} \\
\texttt{{-}{-}{-}{-}{-}{-}{-}} \\
\texttt{* Shifts: 0.00} \\
\texttt{* Max Shifts: 4.00} \\
\texttt{* Min Shifts: 2.00} \\
\texttt{* End Distance Sum: 11.42} \\
\texttt{* Init Distance Sum: 11.42} \\
\texttt{* Expected Distance Sum: 12.57} \\
\texttt{* Penalties: 9.00} \\
\texttt{* Max Penalties: 8.00} \\
\texttt{* Rounds: 2.00} \\
\texttt{* Max Rounds: 12.00} \\
\texttt{* Object Count: 3.00} \\
            }
        }
    }
    & & \\ \\

\end{supertabular}
}

\end{document}
